% Two distinct paths share the same candidate: training updates and use on current information.
% The action feedback path is present only for an interactive task.
\begin{tikzpicture}[
  figure picture,x=1mm,y=1mm,font=\figurefont\small,
  learning object/.style={figure node,font=\figurefont\small,
    minimum height=16mm,inner xsep=2mm,inner ysep=2mm},
  learning flow/.style={figure flow,line width=.9pt},
  learning update/.style={learning flow,draw=FigureLoss,dashed},
  learning label/.style={font=\figurefont\small,text=FigureInk,
    inner sep=1pt,fill=white,align=center}
]
  \path[use as bounding box] (0,0) rectangle (169,76);

  \node[learning object,fill=FigureInputFill,minimum width=24mm,
    minimum height=22mm] (environment) at (18,55)
    {环境 $\mathcal E$\\静态分布 $P$\\交互状态 $S_t$};
  \node[learning object,fill=FigureInputFill,minimum width=24mm]
    (data) at (47,55) {训练记录\\$D$};
  \node[learning object,fill=FigureLossFill,minimum width=26mm]
    (learning) at (80,55) {学习更新\\$A(\ell,\mathcal C)$};
  \node[learning object,fill=FigureModelFill,minimum width=30mm]
    (candidate) at (116,55) {模型或策略\\$h^{(k)}$};
  \node[learning object,fill=FigureOutputFill,minimum width=26mm]
    (prediction) at (155,55) {预测或表示};

  \node[learning object,fill=FigureInputFill,minimum width=28mm]
    (information) at (80,27) {当前信息\\$\bm x$ 或 $H_t$};
  \node[learning object,fill=FigureOutputFill,minimum width=26mm]
    (action) at (155,27) {行动\\$A_t$};

  \draw[learning flow] (environment.east)--(data.west);
  \node[learning label] at (32.5,70) {采样与记录};
  \draw[learning flow] (data.east)--(learning.west);
  \draw[learning update] (learning.east)--(candidate.west);
  \node[learning label] at (98,69) {更新 $k$};
  \draw[learning flow] (candidate.east)--(prediction.west);

  % Current information can be used while the candidate remains fixed.
  \draw[learning flow,rounded corners=1.5mm]
    (environment.south)--(18,27)--(information.west);
  \node[learning label] at (43,31) {观测};
  \draw[learning flow,rounded corners=1.5mm]
    (information.east)--(109,27)--(109,47);
  \draw[learning flow,rounded corners=1.5mm]
    (123,47)--(123,27)--(action.west);

  % No learning step is required between consecutive environment actions.
  \draw[learning flow,rounded corners=1.5mm]
    (action.south)--(155,7)--(2,7)--(2,55)--(environment.west);
  \node[learning label] at (85,12)
    {仅交互：环境步 $t$，$S_t\to S_{t+1}$ 与反馈};
\end{tikzpicture}
